\documentclass[a4paper]{article}
\usepackage[utf8]{inputenc}
\usepackage[T1]{fontenc}
\usepackage[italian]{babel}
\usepackage[margin=2cm]{geometry}
\usepackage{tabularx}
\usepackage{array}
\usepackage{enumitem}

\setlist[itemize]{leftmargin=*, nosep, topsep=2pt}

\begin{document}

\begin{center}
{\Large \textbf{Analisi SWOT -- Sito per ricette}}
\end{center}

\vspace{0.5cm}

\renewcommand{\arraystretch}{1.3}
\begin{tabularx}{\textwidth}{|>{\centering\arraybackslash}m{2.4cm}|X|X|}
\hline
 & \centering\textbf{Fattori positivi} & \centering\arraybackslash\textbf{Fattori negativi} \\
\hline

\textbf{Origine interna}
&
\textbf{Punti di forza (Strengths)}
\begin{itemize}
  \item Interfaccia semplice e intuitiva
  \item Ricerca per ingredienti e filtro per tipo di cibo
  \item Ricette consultabili senza registrazione
  \item Sito multilingua
  \item Recensioni degli utenti sulle ricette
\end{itemize}
&
\textbf{Punti di debolezza (Weaknesses)}
\begin{itemize}
  \item Poche ricette all'inizio, da inserire a mano
  \item Sito nuovo, senza notorietà
\end{itemize}
\\
\hline

\textbf{Origine esterna}
&
\textbf{Opportunità (Opportunities)}
\begin{itemize}
  \item Interesse crescente per la cucina casalinga
  \item Cercare per ingredienti aiuta a ridurre gli sprechi
  \item Possibili estensioni: preferiti, commenti, filtri per dieta e allergie, app
  \item Utenti stranieri grazie al multilingua
\end{itemize}
&
\textbf{Minacce (Threats)}
\begin{itemize}
  \item Concorrenti molto noti (es. GialloZafferano)
  \item Molti siti e app gratuiti di ricette
  \item Protezione dei dati e delle password degli utenti (GDPR)
  \item Diritti d'autore su testi e foto delle ricette
  \item Recensioni false o spam
\end{itemize}
\\
\hline
\end{tabularx}

\end{document}